\documentclass{article}
\usepackage[polish]{babel}
\usepackage[utf8]{inputenc}
\usepackage{polski}
\frenchspacing
\setcounter{tocdepth}{2}
\usepackage{graphicx}
\graphicspath{ {images/} }
\usepackage{float}
\usepackage{listings}
\usepackage{hyperref}

\usepackage{booktabs}% http://ctan.org/pkg/booktabs
\newcommand{\tabitem}{~~\llap{\textbullet}~~}

\linespread{1.3}
\frenchspacing

\setcounter{tocdepth}{3}

\begin{document}
	
	% -------------------------- TITLEPAGE -------------------------- %
	
	\begin{titlepage}
		
		\newcommand{\HRule}{\rule{\linewidth}{0.5mm}}
		
		\begin{center}
			
			% university
			\textsc{\LARGE Politechnika Warszawska}\\[0.5cm]
			\textsc{\Large Wydział Matematyki i Nauk Informacyjnych}\\[1cm]
			
			% logo
			\includegraphics[width=2cm, height=2cm]{logo}\\[1cm]
			
			
			\textsc{\Huge Wstęp do algorytmów ewolucyjnych}\\[0.5cm]
			
			%----------------------------------------------------------------------------------------
			
			\HRule \\[0.4cm]
			{ \LARGE \bfseries Wielostartowy algorytm wspinaczkowy dla ciągłych zadań optymalizacji o różnych wymiarowościach}\\[0.2cm]
			
			%----------------------------------------------------------------------------------------
			
			\HRule \\[0.4cm]
			{  \bfseries Specyfikacja wstępna}\\[2.5cm]
			
			% author
			\begin{flushright}
				\Large \emph{Autorzy:}\\[0.5cm]
				Anna \textsc{Zawadzka}\\
				Piotr \textsc{Waszkiewicz}\\[1.5cm]
			\end{flushright} 
			
			% date
			\vfill
			{\large \today}\\[1cm]
			
		\end{center}
		
	\end{titlepage}
	
	\newpage
	%----------------------------------------------------------------------------------------
	\section{Opis problemu}
	
	Celem projektu jest zbadanie zachowania wielostartowego algorytmu wspinaczkowego dla ciągłych zadań optymalizacji o różnych wymiarowościach z uwzględnieniem rożnych strategii losowania punktów startowych: losowanie z rozkładem równomiernym, przeszukiwanie po hipersiatce, poisson-disc. Testy przeprowadzone będą na benchmarku CEC 2013.
	
	\section{Metoda realizacji zadania}
	
	\subsection{Algorytm wspinaczkowy}
	
	W projekcie wykorzystany zostanie algorytm wspinaczkowy w wersji z wyborem następnika na podstawie sąsiadów najlepszego znalezionego dotychczas punktu.
	
	\begin{lstlisting}[mathescape][language=R]
	x $\leftarrow$ $x_{0}$
	H $\leftarrow$ {$x_{0}$}
	while !stop
	y $\leftarrow$ randomNeighbor(x, $\delta$)
	if cec2013(i, y) > cec2013(i, x)
	x $\leftarrow$ y
	H $\leftarrow$ H $\cup$ {y}
	\end{lstlisting}
	Zbiór H zawiera wszystkie punkty powstałe podczas działania algorytmu i jest nazywany \textit{śladem algorytmu}. Przydaje się podczas graficznej reprezentacji rozłożenia populacji w przestrzeni. W ramach projektu nie będzie on wykorzystywany ze względu na wielowymiarowość problemów, często niemożliwą do narysowania. \\
	Funkcja \textit{randomNeighbor(x, $\delta$)} zwraca losowy punkt w sąsiedztwie $x$-a, znajdujący się w odległości nie większej niż $\delta$ od $x$-a. Do obliczania odległości stosowana będzie metryka euklidesowa. \\
	Funkcja \textit{cec2013(i, z)} dostępna jest w ramach pakietu CEC2013 i zwraca wartość funkcji $i$ w punkcie $z$. Traktowana jest jako funkcja celu.
	\newpage
	Podstawowymi problemami, jakie można napotkać podczas działania algorytmu są:
	\begin{itemize}
		\item Lokalne minima
		\item Plateaux, czyli równiny
		\item Wąskie grzbiety
	\end{itemize}
	
	W niektórych problemach mogą pomóc wielokrotne starty z przypadkowych punktów, stąd cel tego zadania - sprawdzenie które podejście może zminimalizować ryzyko nieodnalezienia poszukiwanego globalnego ekstremum.
	
	
	\subsection{Strategie losowania punktów startowych}
	Punkty startowe wybierane będą przy pomocy następujących metod:
	
	\subsubsection{Losowanie z rozkładem równomiernym}
	Rozkład równomierny gwarantuje wybranie punktu z określonej n-wymiarowej przestrzeni z jednakowym prawdopodobieństwem. Środowisko R dostarcza funkcję \textit{runif(n, a, b)}, która zwraca n liczb z przedziału [a,b] z rozkładu równomiernego, które traktowane będą jako kolejne współrzędne generowanego punktu.
	
	\subsubsection{Przeszukiwanie po hipersiatce}
	Metoda ta zakłada istnienie n-wymiarowej siatki. Przecięcia linii tejże siatki wyznaczają kolejne punkty, będące punktami startowymi algorytmu.
	
	\subsubsection{Poisson Disc Sampling}
	Technika Poisson Disc generuje zbiór punktów o zadanej liczności w n-wymiarowej przestrzeni, ściśle upakowanych, lecz odległych od siebie co najmniej o podaną minimalną odległość. Algorytm zrealizowany zostanie zgodny z poniższym dokumentem:\\
	\url{http://www.cs.ubc.ca/~rbridson/docs/bridson-siggraph07-poissondisk.pdf}
	
	\section{Koncepcja przeprowadzenia eksperymentu}
	
	Pierwszym krokiem doświadczenia będzie wylosowanie N punktów startowych przy pomocy trzech opisanych powyżej metod. Punkty te będą badane dla każdej z funkcji celu z pakietu CEC2013.
	
	Następnym krokiem będzie uruchomienie algorytmu wspinaczkowego dla każdego z początkowych punktów startowych. Co 1000 kroków programu (czyli co 1000 losowań nowego punktu sąsiadującego z najlepszym dotychczas znalezionym rozwiązaniem) wartość funkcji celu dla najlepszego osobnika z populacji będzie notowana. Zakończenie obliczeń jest planowane po uzyskaniu 1000 wyników w trakcie działania programu. Następnie nastąpi wyliczenie wartości średnich dla elementów w wektorze na podstawie wyników otrzymanych dla każdego z punktów startowych. Obliczenia powtórzone będą dla każdej z trzech metod wybierania punktów startowych a otrzymane wyniki zestawione, porównane i opisane. Na koniec zaprezentowane zostaną wyciągnięte wnioski.
	
\end{document}